\begin{insight}{通解是“一个位置加全部方向”}
对$x+y=1$，向量$(1,0)\T$给出一个位置，齐次解方向$(1,-1)\T$决定整条直线的走向。平移位置不会改变方向空间。原点到这条解直线最近的点，恰与齐次方向正交，因此最小范数问题也能直接从图形得到条件。
\end{insight}
\begin{figure}[H]\centering
\begin{tikzpicture}[scale=2.5]
\draw[laaxis](-.55,0)--(1.65,0)node[right]{$x$};\draw[laaxis](0,-.55)--(0,1.65)node[above]{$y$};
\draw[Teal,thick](-.4,1.4)--(1.4,-.4)node[below right]{$x+y=1$};
\draw[dashed,Gold](-.5,.5)--(.65,-.65)node[below]{$x+y=0$};
\draw[lavec,Ink](0,0)--(.5,.5)node[midway,left]{$\vx_{\min}$};
\draw[lavec,Gold](.5,.5)--(.9,.1)node[midway,above right]{$\ker A$方向};
\node[lapt]at(.5,.5){};\node[lanote,anchor=south west]at(.55,.55){$(\frac{1}{2},\frac{1}{2})$};
\draw[Ink!70](.43,.57)--(.5,.64)--(.57,.57);
\node[lapt]at(1,0){};\node[lanote,anchor=north west]at(1.02,-.02){特解$(1,0)$};
\end{tikzpicture}
\caption{非齐次解直线平行于核；最小范数解是原点到解直线的垂足。}\label{fig:affine-minnorm}
\end{figure}
\begin{clevertrick}{不先写所有参数，也能找最小范数解}
对约束$\va\T\vx=b$且$\va\ne0$，先在法向方向寻找解。设$x_0=c\va$，代入$c\va\T\va=b$，因$\va\T\va>0$得
\[x_0=\dfrac{b}{\va\T\va}\va.\]
为什么它最优？任意解$\vx$与$x_0$之差$h$满足$\va\T h=0$，所以$x_0\T h=0$。勾股分解给
\[\norm\vx^2=\norm{x_0+h}^2=\norm{x_0}^2+\norm h^2\ge\norm{x_0}^2,\]
等号仅在$h=0$时成立，故$x_0$就是唯一最小范数解。对$x+2y+2z=3$，$\va=(1,2,2)\T$，$\va\T\va=1+4+4=9$，所以$x_0=(\frac{1}{3},\frac{2}{3},\frac{2}{3})\T$；约束左端为$\frac{1}{3}+\frac{4}{3}+\frac{4}{3}=3$，长度平方$\frac{1}{9}+\frac{4}{9}+\frac{4}{9}=1$。另一法用Cauchy：$|b|\le\norm\va\norm\vx$给长度下界$\frac{|3|}{3}=1$，而$x_0$确能取到。
\textbf{易错提示：}垂足相对于所选内积而言；带权平方和问题的法向条件要相应改变。
\end{clevertrick}
